\chapter{Food: agriculture, logistics, security}
\label{vol12:ch05}

\epigraph{The fact that you have eaten today is the result of decisions made by a quarter-million engineers and several billion farmers, none of whom you have met.}{Author}

\chapmeta{Half-life: medium-life. Archetypes: transport, balance, scaling. Exercise target: 25. Project track: simulation (trace one food product through the system). Hazard class: medium (chronic food-security stakes; ecological footprint). Reader path default: standard, with the agricultural-system and food-security sections core and the alternative-protein section mastery. Named cases: 2007-2008 global food-price crisis (\cite{paper:headey-fan-2008}); US Dust Bowl 1930s (\cite{hist:worster-dustbowl-1979}); Panama disease in bananas (\cite{paper:ploetz-foc-tr4-2015}).}

\noindent
Food is the third master commodity of civilization, after energy and water. The contemporary global food system feeds approximately $8 \times 10^9$ people (with $0.7 \times 10^9$ undernourished as of 2024 per FAO); produces approximately $4 \times 10^9$ tonnes of cereals, $0.4 \times 10^9$ tonnes of meat, and a vast portfolio of fruits, vegetables, oils, and processed foods; consumes approximately $30\%$ of global energy across the field-to-fork chain; produces approximately $25\%$ of global anthropogenic greenhouse gas emissions; occupies approximately $50\%$ of habitable land (including grazing); and uses $70\%$ of global freshwater. The food system is the largest single human intervention in the planet's ecosystems. The chapter takes the agricultural production system as the centre and tracks outward through the processing, distribution, and consumption chains.

\section{The agricultural system as engineered domain}\pathtag{core}

Agriculture is the application of engineering principles to biological production. The Green Revolution of the 1950s-1970s combined new crop varieties (the dwarf wheat varieties of Norman Borlaug; the IR8 rice variety of the International Rice Research Institute), synthetic fertilisers (from the Haber-Bosch process at industrial scale), irrigation expansion (from new dam and pump construction), pesticides, and mechanisation to roughly triple per-hectare cereal yields in the developing world between 1960 and 2000. The result: global cereal production grew from $\approx 0.9 \times 10^9$ tonnes in 1960 to $\approx 3 \times 10^9$ tonnes in 2020, while global population grew from $\approx 3 \times 10^9$ to $\approx 7.8 \times 10^9$. The Malthusian famine projected for the 1970s did not occur; the engineering of agricultural production scaled with population.

The contemporary agricultural-engineering profession works on the next phase of the productivity question. Per-hectare yields in most developed-world systems are at or near the agronomic ceiling for the dominant crops under current climate and water conditions. The productivity growth in the next $50$ years must come from: precision agriculture (variable-rate application; satellite and drone monitoring; automated equipment); biological improvements (CRISPR-based crop modification; nitrogen-fixing cereals; perennial grain crops); resource-efficiency improvements (water-saving irrigation; lower-input agronomic systems); soil-and-system improvements (cover cropping; reduced tillage; integrated crop-livestock systems); and substantial reductions in the food-system's environmental and climate footprint.

\begin{archetype}[Scaling: from field to plate]
A calorie consumed at the plate is the end of a chain that began with photosynthesis in a field, was harvested, transported, processed, packaged, distributed, retailed, and finally cooked. The chain typically consumes $5$-$10$ times the calorific content of the food in fossil energy along the way, with the largest single contribution from the field-stage inputs (synthetic fertiliser energy embodied; on-farm mechanisation; irrigation pumping) and substantial contributions from refrigeration, processing, and transport. The engineering analysis is the energy accounting and the carbon accounting along the entire chain; the engineering optimisation is across the chain rather than within any single stage.
\end{archetype}

\begin{estimation}
\textbf{Question.} Estimate the synthetic nitrogen fertiliser produced globally each year, and the natural-gas demand required to produce it.

\textbf{Reasoning.} Global synthetic nitrogen production: $\approx 120$ million tonnes N/yr (FAO; current as of 2023). Energy intensity of Haber-Bosch ammonia synthesis: $\approx 35\,\text{GJ/tonne NH}_3$, with $14/17 \approx 82\%$ of the mass of ammonia being nitrogen, so $\approx 43\,\text{GJ/tonne N}$. Total energy: $120 \times 10^6 \times 43 \times 10^9 \approx 5 \times 10^{18}\,\text{J} = 5\,\text{EJ/yr}$. Natural gas equivalent at $40\,\text{MJ/kg}$: $\approx 1.3 \times 10^{11}\,\text{kg} = 130 \times 10^6\,\text{tonnes}$. Conclusion: synthetic nitrogen production alone consumes $\approx 1\%$ of global primary energy and is the largest single industrial use of natural gas. Decarbonising agriculture requires decarbonising nitrogen fertiliser, currently via green-hydrogen ammonia synthesis.
\end{estimation}

\section{Soil, water, fertilisers, pesticides}\pathtag{standard}

The soil-water-fertiliser-pesticide complex is the engineering domain of agronomy. Soil is a multi-component material (mineral matter, organic matter, water, air, biota) whose engineering relevance is its capacity to support plant production, to store and release water, to store and supply nutrients, and to resist erosion. Soil degradation (compaction, erosion, salinisation, organic-matter loss, contamination) is the chronic failure mode of agriculture; the US Dust Bowl of the 1930s (\cite{hist:worster-dustbowl-1979}) is the historical reference case for what happens when the soil-engineering principles are violated at landscape scale. The contemporary engineering response: conservation tillage, cover cropping, integrated nutrient management, soil-organic-matter restoration.

Water in agriculture is principally irrigation; the basin-scale dynamics are discussed in Chapter 3 of this volume. The fertiliser system delivers nitrogen, phosphorus, and potassium plus micronutrients in forms that crops can take up; the efficiency of uptake (the fraction of applied nutrient that ends up in the harvested crop) ranges from $30\%$ to $60\%$ depending on crop, soil, climate, and management; the unused fraction becomes runoff (eutrophication of downstream waters), leaching (groundwater nitrate), or atmospheric (nitrous oxide as greenhouse gas; ammonia as air pollution). The pesticide system delivers herbicides, insecticides, fungicides; the principal contemporary engineering and policy issues are pollinator effects (neonicotinoids), human-health effects (organophosphates, glyphosate dispute), and resistance evolution (herbicide-resistant weeds; insecticide-resistant pests).

\begin{principle}[Field-to-watershed continuity]
The agricultural system's inputs (water, nutrients, pesticides) leave the field as outputs (runoff, leaching, drift) at rates determined by the agronomic management and the field-edge conditions. The downstream water and air quality are functions of the upstream agricultural choices. The engineering response to the eutrophication of downstream waters, the groundwater nitrate contamination, the air-pollution from agricultural ammonia, the pesticide residues in surface and ground waters, is at the field-management scale rather than at the downstream treatment scale. The historic engineering tradition that treated the field and the watershed as separate problems produced the failure modes; the contemporary tradition treats them as continuous.
\end{principle}

\section{Mechanisation and precision agriculture}\pathtag{standard}

Agricultural mechanisation followed the same arc as industrial mechanisation: from animal-and-human power to steam (briefly), to internal-combustion (the dominant 20th-century pattern), to increasingly electronic and increasingly autonomous equipment in the 21st century. The contemporary precision agriculture system combines GPS-guided equipment (sub-centimetre accuracy at field scale), variable-rate application (fertiliser and pesticide dosing varied by within-field zones), satellite and drone monitoring (multispectral imaging of crop status), and increasingly autonomous machinery (driverless tractors; robotic weeders; targeted spot-spraying). The economic case: input-cost reduction of $10\%$-$30\%$ for fertiliser and pesticide; yield improvements of $5\%$-$15\%$; reduction in environmental footprint by reducing over-application.

The labour implications are substantial. A US corn-and-soy farm of $1{,}000$ ha in 2024 is typically operated by a single full-time farmer plus seasonal labour; the same farm in 1950 would have required $\approx 20$ farm workers. The transition continues: the projection for 2050 is that automation will reduce farm labour further in developed-world commodity production while remaining substantially smaller in horticulture and high-value crops where the manual operations are harder to automate. The developing-world transition is at an earlier stage and is the largest single labour-market transition currently underway in many countries.

\section{Food processing, cold chain, logistics}\pathtag{standard}

The food-processing system converts raw agricultural products to consumable foods. Engineering scope: unit operations (milling, baking, fermentation, extraction, drying, freezing, packaging); plant design (the processing plant as a process-engineering object); quality assurance (HACCP food-safety systems); and supply-chain integration. The processing sector accounts for approximately $5\%$ of value-added in most developed economies and a comparable share of manufacturing employment; the engineering profession in food processing overlaps substantially with chemical engineering, mechanical engineering, and packaging engineering.

The cold chain (refrigerated storage and transport) is the engineering enabler of perishable-food distribution. A typical perishable supply chain (fresh produce; meat; dairy; pharmaceuticals) requires temperature control from the field or processing plant through wholesale, distribution, retail, and consumer storage. The IIR estimates that the cold chain accounts for approximately $4\%$ of global emissions and is itself growing rapidly as developing-world cold-chain capacity expands. The decarbonisation challenge: refrigerants (with replacement of high-GWP HFCs by lower-GWP HFOs and natural refrigerants under the Kigali Amendment to the Montreal Protocol; \cite{law:kigali-amendment-2016}); energy efficiency of refrigeration equipment; transport decarbonisation (refrigerated trucks and ships).

Food logistics combines the modal-transport considerations of Chapter 4 with the temperature-control and food-safety considerations of the food sector. The just-in-time global food supply chain that brings off-season fresh produce from $10{,}000\,\text{km}$ away to a typical developed-world supermarket is a remarkable engineering achievement and a source of substantial carbon and resilience concerns. The 2020-2021 COVID-19 disruption and the 2022 supply-chain stresses revealed how concentrated and fragile the global food logistics system has become.

\section{Food security and supply-chain resilience}\pathtag{core}

Food security is the condition that all people, at all times, have physical and economic access to sufficient, safe, and nutritious food. The current state: per FAO, approximately $9\%$ of the global population is chronically undernourished as of 2024; the figure had been falling steadily from the 1990s through 2014 but has plateaued and modestly risen since, with COVID-19 and climate-related disruptions as proximate causes. Acute food insecurity (the population requiring emergency food assistance) ranges from $\approx 100$ million to $\approx 300$ million depending on year and conflict intensity.

The 2007-2008 global food-price crisis \cite{paper:headey-fan-2008} is the contemporary reference for systemic food-system failure: cereal prices roughly doubled over $18$ months, driven by a combination of biofuel mandates diverting maize to ethanol, drought in major producing regions (Australia in particular), high oil prices feeding through to fertiliser and transport costs, low global stocks-to-use ratios, and export restrictions by major producers (Russia, Ukraine, India, Argentina, Vietnam). The crisis produced food-related unrest in approximately $30$ countries and contributed to the political conditions of the Arab Spring. The engineering response: the buffer-stock policies, the production-capacity investments, and the trade-system arrangements that reduce the vulnerability of low-income food-importing populations to price shocks.

The supply-chain-resilience question generalises beyond crises. The contemporary developed-world food system carries roughly $5$-$15$ days of inventory at any point in the chain; the developing-world subsistence system carries less. A multi-month disruption (such as a major regional pandemic; a regional conflict; a climate-related crop failure across a major production region) would produce cascading effects that the engineering profession is asked to anticipate and design against.

\begin{archetype}[Ethics: feeding $10$ billion]
The global population is projected to peak at $\approx 10$ billion around 2080. Feeding $10$ billion at a developed-world dietary standard with current technology would require approximately doubling current agricultural production, with corresponding land, water, and emissions implications. The engineering and ethical question: how to deliver adequate nutrition for the global population while respecting the planetary boundaries that Chapter 12 of this volume describes. The answer combines productivity improvements, dietary transitions (toward lower-meat-and-dairy diets), food-loss-and-waste reductions (currently $\approx 30\%$ of food is lost or wasted), and the alternative-protein technologies that the next section discusses. The engineering profession's contribution: identify the technically feasible pathways; the choice among them is a societal choice that the profession informs but does not determine.
\end{archetype}

\section{Alternative protein and emerging food systems}\pathtag{mastery}

The alternative-protein category includes plant-based meat substitutes (Beyond Meat, Impossible Foods), cultivated meat (cell-cultured meat in bioreactors), insect protein, microbial protein (such as Quorn mycoprotein), and algal protein. The engineering case: conventional meat production has substantially higher land, water, and emissions footprints per kilogram of protein than these alternatives, and substituting even part of meat consumption with alternatives reduces the environmental footprint of the food system. The current state: plant-based alternatives are commercial at scale and produce roughly $25\%$ of the emissions of equivalent meat; cultivated meat is at early commercial stage with cost-and-scale challenges; microbial protein is established in some markets (mycoprotein in Europe) and growing.

The cellular-agriculture engineering question is principally bioreactor scale-up and cost reduction. Current cultivated-meat production is at small-pilot scale (kg-per-day) with production costs an order of magnitude above conventional meat. The scaling challenge is the bioreactor design (oxygen and nutrient transfer in large vessels), the growth-medium cost (the historical use of fetal bovine serum is being replaced by chemically-defined media), and the scaffold-and-structure engineering for whole-cut products. The economic projections suggest that price parity with premium meat could be achievable by the early 2030s for some products; mass-market parity is further off.

\section{Failure: food-system failures}\pathtag{core}

The chapter's failure cases each illustrate a distinct failure mode of food engineering at civilization scale.

The Irish Potato Famine of 1845-1849 is the classical case of agricultural-monoculture failure: the Irish rural population's caloric dependence on a single potato variety ('Irish Lumper') made the population catastrophically vulnerable to a single pathogen (\emph{Phytophthora infestans}, potato late blight). The proximate cause was the pathogen; the structural cause was the monoculture; the institutional cause was the political-economic regime that continued grain exports during the famine. The engineering lesson: agricultural biodiversity is a resilience property of the food system, not a luxury.

The US Dust Bowl of the 1930s \cite{hist:worster-dustbowl-1979} is the case of soil-engineering failure at landscape scale: the conversion of native prairie to wheat agriculture, with deep ploughing and continuous cropping, depleted soil organic matter and destroyed the soil's wind-resistant structure; the multi-year drought of the early 1930s then mobilised the topsoil into the dust storms that produced the catastrophic erosion. The institutional response (the Soil Conservation Service, established 1935; the conservation programmes; the long-term restoration agronomy) was the engineering and policy infrastructure created in response to the failure.

The Panama disease (Fusarium wilt; \emph{Fusarium oxysporum} f. sp. \emph{cubense}) in bananas \cite{paper:ploetz-foc-tr4-2015} is the contemporary monoculture case: the global Cavendish banana export industry depends on a single clone that is vulnerable to a single soil pathogen (TR4 strain); the pathogen has spread from Asia to Africa to the Americas, and the displacement of the Cavendish from major production regions is foreseeable. The engineering response combines disease-management agronomy (sanitation; resistance-genetic alternatives), accelerated breeding programmes for resistant cultivars, and a transition to a more diverse banana-supply portfolio. The same engineering lesson as the Irish potato famine, $180$ years later, in a different crop, with the same structural condition.

The 2007-2008 food-price crisis \cite{paper:headey-fan-2008} is the case of integrated-system failure: the system's normal operation is robust, but the combination of multiple simultaneous stresses exceeded its absorption capacity. The engineering response is the global food-security infrastructure that has improved since (better stocks-to-use ratios in major producers; the AMIS information system to reduce market uncertainty; the World Food Programme's emergency-response capacity), with continuing vulnerability in the regions and populations that were exposed in 2008 and remain so.

\begin{principle}[Diversity as engineered resilience]
The food-system failures share a structural pattern: high concentration in a single variety, a single pathogen-vulnerability, a single production region, or a single supply route produces vulnerability whose magnitude is hidden during normal operation and revealed only by the perturbation. The engineering response is diversity: multiple varieties; multiple supply regions; multiple supply routes; multiple alternative protein sources. Diversity is engineered into the food system at cost; the cost is the insurance premium against the failure mode the chapter's cases illustrate.
\end{principle}

\section{Project}\pathtag{core}

\begin{project}[Food-system trace]
Choose a single food product (a loaf of bread; a beef steak; a tomato; a bottle of olive oil). Trace the product from the field or production unit to your plate, documenting each engineering touch-point along the chain. Document the inputs (seed; fertiliser; water; energy; labour); the processing operations; the packaging; the distribution; the retail; the consumer preparation. Estimate the carbon footprint at each stage. Identify the engineering decisions that determine the footprint; identify the alternatives that would reduce it; identify the system-resilience properties (single-source vs diversified supply; perishability; substitutability).
\end{project}

\section{Exercises}

\subsection*{Calculation}\pathtag{core}

\begin{exercise}
A maize field in the US Midwest yields $12\,\text{t/ha}$ at $3.7\,\text{kcal/g}$ dry matter and $14\%$ moisture content. Compute the total kcal yield per hectare and the per-hectare population that can be fed at a $2{,}500\,\text{kcal/person/day}$ diet, ignoring further losses.
\end{exercise}

\begin{exercise}
A beef-cattle operation produces $1\,\text{kg}$ of edible beef from approximately $7\,\text{kg}$ of grain (for grain-fed beef). Compute the land-area implication of a $25\%$ shift from cereal-based to beef-based caloric intake for a population of $300$ million.
\end{exercise}

\begin{exercise}
A nitrogen-fertiliser application rate of $200\,\text{kg N/ha}$ at $40\%$ nitrogen-use efficiency releases $120\,\text{kg N/ha}$ to the environment, of which approximately $1\%$ becomes N$_2$O (the IPCC default emission factor). Compute the CO$_2$-equivalent emissions per hectare at the IPCC global-warming-potential value of $273$ for N$_2$O.
\end{exercise}

\begin{exercise}
A refrigerated container shipping fresh strawberries from Spain to Germany maintains $2\degC$ for $48$ hours over a $1{,}500\,\text{km}$ road haul. Compute the refrigeration energy required and the additional emissions beyond the transport-only emissions.
\end{exercise}

\begin{exercise}
A food-processing plant uses $100\,\text{tonnes}$ of input maize/day at $14\%$ moisture and produces $80\,\text{tonnes}$/day of corn starch at $10\%$ moisture plus by-product streams. Compute the water removed per day and the energy required at $2.5\,\text{MJ/kg water}$ for drying.
\end{exercise}

\subsection*{Derivation}\pathtag{standard}

\begin{exercise}
Derive the relationship between agricultural-yield improvement rate, population growth rate, and arable-land requirement over a $50$-year horizon. Identify the conditions under which arable-land expansion is required.
\end{exercise}

\begin{exercise}
Derive the food-system carbon footprint as a sum of stage contributions (field, processing, transport, retail, consumer). Identify the dominant stages for a high-meat versus a low-meat diet.
\end{exercise}

\begin{exercise}
Derive the optimal stocks-to-use ratio for a strategic grain reserve, given storage cost, opportunity cost, and the distribution of supply-disruption events.
\end{exercise}

\begin{exercise}
Derive the substitution elasticity between conventional and alternative proteins at the consumer level, given price and preference parameters. Identify the price point at which substitution becomes substantial.
\end{exercise}

\subsection*{Estimation}\pathtag{standard}

\begin{exercise}
Estimate the global land area required for current meat consumption at developed-world per-capita rates extended to the projected $10$-billion population. Compare to current agricultural land and to ice-free land area.
\end{exercise}

\begin{exercise}
Estimate the global potential for food-loss-and-waste reduction at each stage (production, post-harvest handling, processing, distribution, consumption) and the implied caloric and land-area savings.
\end{exercise}

\begin{exercise}
Estimate the cold-chain energy and emissions for a typical developed-world supermarket and project the global emissions implication of cold-chain expansion to developing-world levels.
\end{exercise}

\begin{exercise}
Estimate the maximum sustainable global fish catch under current marine-ecosystem conditions and compare to current catch levels. Identify the engineering and policy implications.
\end{exercise}

\subsection*{Simulation}\pathtag{standard}

\begin{exercise}
Simulate a regional food-supply system under a multi-year climate-driven yield reduction. Sweep the storage-and-import-policy parameters and identify the policies that maintain food-security versus the policies that fail.
\end{exercise}

\begin{exercise}
Simulate a precision-agriculture deployment for a $1{,}000$-ha grain operation. Compare the input-use and yield outcomes against a conventional-management baseline.
\end{exercise}

\begin{exercise}
Simulate the cellular-agriculture cost trajectory under different learning-curve assumptions and identify the years at which price-parity with conventional meat becomes plausible.
\end{exercise}

\subsection*{Design}\pathtag{standard}

\begin{exercise}
Design a strategic grain-reserve system for a major food-importing country of $50$ million people. Specify the reserve size, the storage locations, the rotation policy, the release rules, the cost-recovery model.
\end{exercise}

\begin{exercise}
Design an integrated soil-conservation programme for a $10{,}000$-ha agricultural watershed facing chronic erosion. Specify the agronomic interventions, the structural interventions, the policy and incentive structure, the monitoring and evaluation.
\end{exercise}

\begin{exercise}
Design a low-input agronomic system for a transition from conventional to lower-environmental-footprint production on a $500$-ha mixed farm. Specify the rotation, the nutrient management, the pest management, the economic transition support required.
\end{exercise}

\subsection*{Judgement}\pathtag{mastery}

\begin{exercise}
Argue: the global food system's environmental footprint is so large that dietary transitions toward lower-meat diets are an engineering necessity, not a personal-preference matter. Argue the opposite. Identify the empirical evidence and the value judgments.
\end{exercise}

\begin{exercise}
Discuss the 2007-2008 food-price crisis as a food-system failure. Identify the engineering and policy interventions that have reduced subsequent vulnerability and the residual exposures.
\end{exercise}

\begin{exercise}
Argue: cellular agriculture is the appropriate engineering response to the projected protein-demand expansion; conventional livestock should be largely phased out for environmental and animal-welfare reasons. Argue the opposite. Identify the engineering and ethical considerations.
\end{exercise}

\begin{exercise}
Discuss the engineering profession's role in food-system design in low-income countries. Identify the technical contributions; identify the institutional limits; identify the imperative not to replicate developed-world failures (Green-Revolution overuse of inputs; monoculture lock-in).
\end{exercise}

\begin{exercise}
Discuss the Panama disease in bananas as a contemporary monoculture vulnerability. Identify the engineering response options; identify the structural changes required to make the global banana supply more resilient.
\end{exercise}

\begin{exercise}
Discuss the US Dust Bowl as a soil-engineering failure. Identify the agronomic decisions that produced the failure; identify the institutional response that emerged from it; identify the contemporary regions whose agronomic practices put them at comparable risk.
\end{exercise}
